\chapter{Propuesta de optimización sostenible}

Este capítulo desarrollará la propuesta arquitectónica optimizada a partir del diagnóstico del escenario base y de la evaluación de estrategias pasivas y materiales sostenibles. La propuesta deberá responder a las condiciones climáticas, normativas y ambientales de la isla San Cristóbal.

\section{Criterios de optimización}

Los criterios de optimización se orientan a reducir la demanda energética, mejorar el confort térmico, fortalecer la ventilación natural, controlar la radiación solar, mejorar la eficiencia lumínica y disminuir el impacto ambiental de los materiales.

\section{Estrategias pasivas propuestas}

Las estrategias pasivas podrán incluir control solar, optimización de aperturas, ventilación cruzada, ventilación nocturna, aumento de inercia térmica, protección de la envolvente, uso de vegetación y mejora de la iluminación natural.

\section{Materiales sostenibles propuestos}

La selección de materiales deberá considerar bajo carbono incorporado, disponibilidad regional, durabilidad, comportamiento térmico, mantenimiento, potencial de reutilización y coherencia con el contexto insular. Se priorizarán soluciones que reduzcan impactos logísticos y ambientales.

\section{Evaluación comparativa de escenarios}

La propuesta será evaluada comparando el escenario base con los escenarios optimizados. Esta comparación permitirá determinar la incidencia de cada estrategia y de la combinación de estrategias pasivas con materiales sostenibles.

\section{Lineamientos replicables}

Con base en los resultados se formularán lineamientos aplicables a futuras edificaciones públicas en Galápagos y contextos insulares similares. Estos lineamientos deberán integrar criterios de diseño pasivo, eficiencia energética, selección de materiales y evaluación mediante simulación.
